%\documentclass[12pt,a4paper]{article}
%\usepackage[T1]{fontenc}
%\usepackage[utf8]{inputenc}
%\usepackage{t1enc}
%\usepackage{amsmath,amssymb}
%\usepackage{pst-plot,pst-tree,pstricks,pst-node,pst-text}
%\usepackage{pst-eucl}
%\usepackage{pstricks-add}
%\usepackage[left=3.5cm, right=3.5cm, top=3cm, bottom=3cm]{geometry}
%\newcommand{\vect}[1]{\overrightarrow{\,\mathstrut#1\,}}
%\def\Oij{$\left(\text{O}~;~\vect{\imath},~\vect{\jmath}\right)$}
%\def\Oijk{$\left(\text{O}~;~\vect{\imath},~\vect{\jmath},~\vect{k}\right)$}
%\def\Ouv{$\left(\text{O}~;~\vect{u},~\vect{v}\right)$}
%% Macros
%\def\Q{\mathbf{Q}}
%\def\Z{\mathbf{Z}}
%\def\C{\mathbf{C}}
%\def\N{\mathbf{N}}
%\def\R{\mathbf{R}}
%\renewcommand{\theenumi}{\arabic{enumi}}
%\def\labelenumi{{\bf \theenumi /}}
%\def\labelenumii{{\theenumii)}}
%\parindent0pt
%\usepackage[french]{babel}
%\DecimalMathComma
%\usepackage[np]{numprint}
%\begin{document}
%
%\begin{center}
%\Large\textsf{CONCOURS GÉNÉRAL DES LYCÉES}\\
%\rule{2cm}{0.1mm}\\[2mm]
%\textsf{SESSION DE 2001}\\
%\rule{2cm}{0.1mm}\\[2mm]
%\textbf{COMPOSITION DE MATHÉMATIQUES}\\
%\normalsize
%\textsf{(Classe de Terminale S)}\\[2mm]
%\textsc{Durée : 5 heures}\\
%\rule{2cm}{0.1mm}
%\end{center}
\documentclass[a4paper,12pt]{article}
\usepackage[T1]{fontenc}
\usepackage[utf8]{inputenc}
\usepackage{t1enc}
\usepackage[dvips]{color}
\usepackage{amsmath,amssymb,makeidx,stmaryrd}
\usepackage[normalem]{ulem}
\usepackage{fancybox}
\definecolor{gris}{gray}{0.9}
\def\bas{\fcolorbox{gris}{gris}{\parbox{0.986\textwidth}{
	\color{red}\footnotesize 2001\hfill Page : \thepage/1}}}
\def\haut{\fcolorbox{gris}{gris}{\parbox{0.986\textwidth}{
	\color{red}
	\bf Concours Général de Mathématiques \hfill 2001
	}}}
\makeatletter
\def\ps@cg{%
\let\@oddhead\haut%
\let\@evenhead\haut%
\let\@oddfoot\bas
\let\@evenfoot\bas}
\makeatother
\pagestyle{cg}
\usepackage[left=3.25cm, right=3.25cm, top=3cm, bottom=3cm]{geometry}
\usepackage{pst-plot,pst-tree,pstricks,pst-node,pst-text}
\newcommand{\vect}[1]{\overrightarrow{\,\mathstrut#1\,}}
\newcommand{\vectt}[1]{\overrightarrow{\,\mathstrut\text{#1}\,}}
%\renewcommand{\theenumi}{\textbf{\arabic{enumi}}}
%\renewcommand{\labelenumi}{\textbf{\theenumi.}}
%\renewcommand{\theenumii}{\textbf{\alph{enumii}}}
%\renewcommand{\labelenumii}{\textbf{\theenumii.}}
\newcommand{\R}{\mathbb{R}}
\newcommand{\N}{\mathbb{N}}
\newcommand{\D}{\mathbb{D}}
\newcommand{\Z}{\mathbb{Z}}
\newcommand{\Q}{\mathbb{Q}}
\newcommand{\C}{\mathbb{C}}
\def\Oij{$\left(\text{O}~;~\vect{\imath},~\vect{\jmath}\right)$}
\def\Oijk{$\left(\text{O}~;~\vect{\imath},~\vect{\jmath},~\vect{k}\right)$}
\def\Ouv{$\left(\text{O}~;~\vect{u},~\vect{v}\right)$}
\parindent0pt
\everymath{\displaystyle}
\usepackage[french]{babel}

\DecimalMathComma
\usepackage[np]{numprint}
\begin{document}

\begin{center}\textbf{{\color{blue} Exercice 1}}
\end{center}

\begin{center}
La calculatrice de poche est autorisée.\\
La clarté et la précision seront prises en compte\\
dans l'appréciation des copies.
\end{center}

{\bf
Les premières questions de chacune des quatre parties de ce problème sont
indépendantes des autres parties. Il n'est donc pas nécessaire de
commencer son étude dans l'ordre indiqué.

Les candidats peuvent admettre les résultats d'une question, à condition de
l'indiquer clairement sur la copie.}
\newpage
\vspace{1cm}

On appelle \textbf{trio} tout triplet de nombres réels $(a, b, c)$ non tous
nuls et vérifiant la relation :

\[ab + bc + ca = 0.\]

Lorsque $a + b + c = 1$, on dit que le trio $(a, b, c)$ est un trio \textbf{réduit}.

Les coordonnées sont rapportées à un repère orthonormal direct
\Oijk{} de l'espace.


\section*{\color{blue}Première partie}

On note $\mathcal{C}$ l'ensemble des points de coordonnées $(a, b, c)$ où
$(a, b, c)$ est un trio.

On note $\Gamma$ l'ensemble des points de coordonnées $(a,b,c)$ où
$(a, b, c)$ est un trio réduit.

On note $\mathcal{P}$ le plan d'équation $x+y+z=1$.

\begin{enumerate}
\item Existe-t-il des trios $(a, b, c)$ tels que $a + b + c = 0$ ?
\item Montrer que $\mathcal{C}$ est une réunion de droites passant par $O$
et privées de ce point.
\item Montrer que $\Gamma$ est l'intersection d'un plan et d'une sphère de
centre $O$.  Quelle est la nature géométrique de $\Gamma$ ?
\item Donner la nature géométrique de $\mathcal{C}$ et l'illustrer d'un
croquis.
\item Soit $ L$ un point fixé de $\Gamma$.
Montrer que le volume $V$ du tétraèdre OLL$'$L$''$, où L$'$ et L$''$ sont
deux points distincts de $\Gamma$ et différents de L, est maximal lorsque
les arêtes issues de O sont deux à deux orthogonales et déterminer alors
les coordonnées de $L'$ et $L''$ en fonction de celles de L.
\item Montrer que le produit $abc$ admet un maximum et un minimum lorsque
le point de coordonnées $(a, b, c)$ décrit $\Gamma$.

Préciser les trios réduits réalisant ces extremums.
\end{enumerate}

\newpage

\section*{\color{blue}Deuxième partie}


Dans cette partie et les suivantes, un trio $(a, b, c)$ est dit
\textbf{rationnel} lorsque $a, b$ et $c$ sont des nombres rationnels
(éléments de l'ensemble $\Q$)~; il est dit \textbf{entier}, lorsque $a, b$
et $c$ sont des nombres entiers relatifs (éléments de l'ensemble $\Z$)~;
enfin un trio entier est dit \textbf{primitif} si $a, b$ et $c$ n'admettent
que $1$ et $- 1$ comme diviseurs communs.

\medskip

\begin{enumerate}
\item Déterminer la nature de l'ensemble $H_1$ des points de coordonnées
$(x, y, 1)$ tels que $(x, y, 1)$ soit un trio.  Montrer que le point $\Omega_1$
de coordonnées $(-1,-1,1)$ est un centre de symétrie de $H_1$. Quels sont
les points de $H_1$ à coordonnées entières ?
\item Pour tout entier naturel non nul $h$, on note $Z_h$ l'ensemble des
trios entiers $(a, b, c)$ tels que $c=h$.

Déterminer $Z_h$ pour $h = 1$ et $h = 2$.
\item Montrer que $Z_h$ est un ensemble fini et exprimer le nombre $N(h)$
de ses éléments en fonction de celui des diviseurs de $h^2$ dans $\Z$.

Montrer que $4$ divise $N(h)-2$.
\item Pour tout entier naturel non nul $h$, on note $N'(h)$ le nombre des
trios entiers $(a, b, c)$ tels que l'un au moins des entiers $a,b$ ou $c$
soit égal à $h$?

Exprimer $N'(h)$ en fonction de $N(h)$ selon la parité de
$h$.
\item Montrer qu'à tout trio entier $(a, b, c)$ on peut associer un triplet
$(r,s,t)$ d'entiers tels que $r $ et $s$ soient premiers entre eux, $s$
positif ou nul et tels que l'on ait : $a=r(r+s)t$, $b=s(r+s)t$, $c=-rst$.

Énoncer et démontrer une réciproque. 

Pour quels trios $(a, b, c)$ le triplet
$(r, s, t)$ n'est-il pas unique ?
\item Déterminer les triplets $(r, s, t)$ ainsi associés aux trios primitifs.

En déduire que si $(a, b, c)$ est un trio primitif, alors $|abc|,\ |a + b|,\
|b + c|$, et $|c + a|$ sont des carrés d'entiers.
\item pour tout entier naturel non nul $h$, on note $P(h)$ le nombre de
trios primitifs $(a, b, c)$ tels que $c = h$.

Montrer que $P(h)$ est une puissance de $2$.  Pour quels entiers $h$ a-t-on $P(h) = N(h)$ ? 

Expliciter une suite d'entiers $(h_n)$ tels que la suite $\left(\frac
{P(h_n)}{H(h_n)}\right)$ converge vers zéro.
\item Soit $(a, b, 1)$ un trio.

Montrer qu'il existe deux suites $(x_n)$ et $(y_n)$ convergeant respectivement vers $a$ et $b$ et telles que pour tout $n$, $(x_n, y_n,1)$ soit un trio rationnel.
\item Soit $(a, b, c)$ un trio réduit.  Montrer qu'il existe trois suites
$(x_n),\ (y_n)$ et $(z_n)$ convergeant respectivement vers $a, b$ et $c$ et
telles que pour tout $n$, $(x_n, y_n, z_n)$ soit un trio rationnel réduit.
\end{enumerate}

\newpage

\section*{\color{blue}Troisième partie}

On note j le nombre complexe $\text{e}^{\frac{2\text{i}\pi}{3}}$, c'est à dire j$ = - \dfrac12
+ \text{i}\dfrac{\sqrt 3}{2}$.

Pour tout trio $T = (a, b, c)$ on note $\widehat{T} = (a, c, b),\, S(T) = a + b + c$ et

$z(T) = a + b\text{j} + c\text{j}^2$.

\medskip

\begin{enumerate}
\item Calculer le module de $z(T)$ en fonction de $S(T)$.

Peut-on avoir $z(T) =0$ ?

Calculer le cosinus et le sinus d'un argument $\theta$ de
$z(T)$ en fonction de $a, b$ et $c$.
\item Soit $z_0$ un nombre complexe non nul.  Déterminer les trios
$T=(a, b, c)$ tels que $z(T)=z_0$.
\item Étant donnés deux trios $T_1$ et $T_2$, montrer qu'il existe un
unique trio, noté $T_1*T_2$, vérifiant $S(T_1*T_2) = S(T_1)S(T_2)$ et
$z(T_1*T_2) = z(T_1)z(T_2)$. Calculer $T_1*T_2$ en fonction de $T_1$ et
$T_2$.

Que peut-on dire d'un argument de $z(T_1*T_2)$ ?

Que peut-on dire d'un argument de $z(T_1*\widehat{T_1})$ ?
\item si $T_1$ et $T_2$ sont réduits, en est-il de même de $T_1*T_2$ ?

Si $T_1$ et $T_2$ sont entiers, en est-il de même de $T_1*T_2$ ?  Si $T_1$ et
$T_2$ sont primitifs, en est-il de même de $T_1*T_2$ ?
\item Comparer les trios $T_1*T_2$ et $T_2*T_1,\ (T_1*T_2)*T_3$ et
$(T_1*T_2)*T_3,\ T_1$ et $T_1*(1,0,0)$.
\item Étant donnés les trios $T_1$ et $T_2$, résoudre l'équation
$T_1*T = T_2$ où le trio $T$ est l'inconnue.
\item Étant donné un trio $T$, on définit une suite de trios $(T_n)$ par
$T_0 = (1,0,0)$ et $T_{n+1}=T*T_n$.

Calculer $S(T_n)$.

Étant donné un entier $p$, résoudre l'équation $T_p = T_0$ où le trio $T$ est l'inconnue.
\end{enumerate}

\newpage

\section*{\color{blue}Quatrième partie}
On note $A$ l'ensemble des entiers $m$ non nuls tels qu'il existe
deux entiers $u, v$ tels que $m = u^2 + 3v^2$.

On note $A'$ l'ensemble des nombres complexes $z$ non nuls tels qu'il existe deux entiers
$u, v$ tels que $z = u + \text{i}v\sqrt 3$ (on remarquera que
$|z|^2 = u^2 + 3v^2$).

On note $B$ l'ensemble des entiers $n$ non nuls
tels qu'il existe deux entiers $r,s$ tels que $n=r^2 + rs + s^2$.
\begin{enumerate}
\item Montrer que le produit de deux éléments de $A'$ appartient à $A'$,
puis que le produit de deux éléments de $A$ appartient à $A$.
\item Montrer que si $p$ est en nombre premier élément de $A$, alors $p=3$
ou $3$ divise $p-1$.
\item Montrer que $A=B$ (on pourra notamment remarquer que

$r^2+rs+s^2 = (r+s)^2- (r + s)s + s^2$).
\item Montrer que $4$ divise les éléments pairs de $A$ et que les quotients
appartiennent à $A$, puis que tout élément de $A$ est produit d'un élément
impair de $A$ par une puissance de $4$.
	\begin{enumerate}
		\item Soit s'il en existe, un entier impair $m = u^2 + 3v^2$ tel que les
entiers $u$ et $v$ soient premiers entre eux et qui admet un diviseur
premier $p$ n'appartenant pas à $A$.

Montrer qu'il existe alors un plus petit entier strictement positif $n_0$ tel que $n_0p$ appartienne à $A$.

Montrer que $n_0$ est impair.
		\item Établir l'existence de deux entiers $u'$ et $v'$ inférieurs en valeur
absolue à $\frac p2$ tels que $p$ divise $u'-u$ et $v'-v$.

Montrer que $p$ divise l'entier non nul $u'^2+3v'^2$ et que $n_0<p$.
		\item Établir l'existence de deux entiers non nuls premiers entre eux $u_0$
et $v_0$ tels que $n_0p=y_0^2+3v_0^2$.
		\item Établir l'existence de deux entiers $u_1$ et $v_1$ inférieurs en
valeur absolue à $\frac{n_0}2$ tels que $n_0$ divise $u_1-u_0$ et
$v_1-v_0$. 

Montrer que $n_0$ divise l'entier non nul $u_1^2+3v_1^2$ que
l'on notera $n_0n_1$.
		\item En déduire qu'un tel entier $m$ ne peut pas exister (on pourra
considérer l'entier $n_0^2n_1p$).
	\end{enumerate}
\item Montrer que tout élément de $A$ s'écrit $m = C^2p_1...p_k$, où
$C$ est un entier naturel non nul et les $p_i$ des nombres premiers
distincts éléments de $A$.
	\begin{enumerate}
		\item Soit $p$ un nombre premier tel que $3$ divise $p-1$ et $K$ l'ensemble
des triplets $(x, y, z)$ où les entiers $x, y$ et $z$ sont strictement compris
entre $0$ et $p$, et tels que $p$ divise ($xyz - 1)$. Montrer que $K$
possède $(p-1)^2$ éléments, et que $3$ divise le nombre d'éléments de $K$
ne vérifiant pas $x=y=z$.
		\item En déduire qu'il existe en entier $x$ strictement compris entre $1$
et $p$ tel que $p$ divise $x^2+x+1$, puis que $p$ appartient à $A$.

Décrire les éléments de $A$.
	\end{enumerate}
\item Soit $D$ l'ensemble des entiers $d$ tels qu'il existe un trio entier
$(a, b, c)$ vérifiant $a + b + c= d$ et $abc\not=0$.

Montrer, grâce à la question \textbf{5/} de la deuxième partie, que tout élément de $D$ possède un diviseur premier élément de $A$.

Réciproquement, que peut-on dire d'un entier non nul admettant un diviseur premier élément de $A$ ?
\item En déduire les éléments de $D$ compris aux sens large entre \np{2001} et
\np{2010}.
\end{enumerate}
\end{document}